% Generated draft: requires mathematical and source review.
\documentclass[11pt]{article}
\usepackage[margin=25mm]{geometry}
\usepackage{fontspec}
\setmainfont{Latin Modern Roman}
\usepackage{amsmath,amssymb,mathtools,unicode-math}
\setmathfont{Latin Modern Math}
\usepackage{enumitem,needspace,xurl,hyperref}
\setlength{\parindent}{0pt}
\setlength{\parskip}{5pt}
\setlength{\emergencystretch}{3em}
\newcommand{\N}{\mathbb{N}}
\newcommand{\Z}{\mathbb{Z}}
\newcommand{\Q}{\mathbb{Q}}
\newcommand{\R}{\mathbb{R}}
\newcommand{\C}{\mathbb{C}}
\newcommand{\F}{\mathbb{F}}
\newcommand{\A}{\mathbb{A}}
\newcommand{\PP}{\mathbb{P}}
\newcommand{\E}{\mathbb{E}}
\newcommand{\eps}{\varepsilon}
\newcommand{\dd}{\,\mathrm{d}}
\begin{document}
\section*{UnsolvedMath ID: 30005955}
\textbf{Source problem code:} catalogue entry
\subsection{Definitions and mathematical statement}
The supplied record for permanent UnsolvedMath ID \(30005955\), source code \texttt{catalogue entry}, contains a truncated table-of-contents fragment rather than a mathematical statement. Its initial word ``conjecture'' is followed by a page reference to 1840, without the conjecture's content or attribution. Subsequent entries list Nicolas Addington, jointly with Benjamin Tighe, on the Torelli problem for Calabi--Yau threefolds and the integral cohomology of the Fermat quintic, at page 1843; and Lie Fu, jointly with Salvatore Floccari, on the Hyper-Kummer construction, at page 1846. The fragment also contains the report heading and page number 1791.

No definitions, hypotheses, domains, boundary conditions, quantifiers, or mathematical conclusion identifying an intended problem are supplied. Consequently, no precise mathematical statement can be recovered from this record. The listed talk titles do not determine a conjecture and do not establish authorship of the preceding unidentified conjecture.

The supplied statement status is \texttt{unrecoverable}. The supplied verification note reports that the official report confirms the extraction is a table-of-contents fragment containing no mathematical conjecture; this is a source-record assertion, not an independently verified finding here. The record's \texttt{open} label likewise does not establish the status of any identifiable mathematical claim.
\subsection{Short English statement}
This entry preserves an incomplete table of contents. It names talks but supplies no recoverable mathematical problem. A statement-only draft requires the intended problem's full text and precise location before mathematical formalization or review is possible.
\subsection{Sources}
\begin{itemize}
\item Local UnsolvedMath registry, permanent ID \(30005955\), source code \texttt{catalogue entry}, title ``Calabi-Yau Torelli Table-of-Contents Fragment,'' in the collection ``Oberwolfach Reports Open Problems.'' The supplied record is the immediate evidence for this draft.
\item Supplied report citation: \emph{Algebraic Geometry: Wall Crossing and Moduli Spaces, Varieties and Derived Categories} (2025), \emph{Oberwolfach Reports}, pp.~1787--1854, DOI \texttt{10.4171/owr/2024/32}. Supplied source URL and literature-source link labelled ``original OWR report'': \url{https://doi.org/10.4171/owr/2024/32}. The collection describes these reports as documenting mathematical workshops and their open-problem discussions. No complete problem statement is supplied from this report.
\item Attribution preserved from the fragment: Nicolas Addington, jointly with Benjamin Tighe, ``On the Torelli problem for Calabi--Yau 3-folds, and the integral cohomology of the Fermat quintic,'' p.~1843; Lie Fu, jointly with Salvatore Floccari, ``Hyper-Kummer construction,'' p.~1846. No proposer is identified for the intended problem.
\item Literature assessment dated exactly \texttt{2026-08-22}: supplied dated curation, not independently verified. Its status label is ``open,'' and its assessment states that the truncated table-of-contents fragment does not identify a mathematical claim for verification. It lists the original OWR report as checked and requests independent verification of time-sensitive status and cited claims.
\item The supplied proposed year and year tag are 2024, whereas the supplied bibliographic citation gives 2025; the DOI contains 2024. These details are retained without resolving their relationship.
\end{itemize}
\end{document}
